% Chapter 5: Before the protocol: the legacy-split work (10 Oct to 14 Nov). Draft 23 Nov 2026.
% Every number is traced in ch05_legacy_NOTES.md. Tables are in the appendix and referenced only by \ref{tab:...}.
\chapter{Before the protocol: the legacy-split work (10 Oct to 14 Nov)}\label{ch:legacy}

This chapter reports the work done before the benchmark protocol existed. It is short because its numbers decide nothing in this
thesis, and it is here because its measurement failures are the reasons for the protocol of Chapter~\ref{ch:protocol}.
Table~\ref{tab:T21} lists all 38 arms and decisions of the period, and Table~\ref{tab:T17} the seven registered openings of E1
before the final one. Figure~\ref{fig:F8} shows how the project's working beliefs changed over the period.

\FIG{F8}{What the project believed and what the evidence said, 7 Oct to 21 Nov 2026 (pre-protocol rows from Table~\ref{tab:T21}, E1
openings from Table~\ref{tab:T17}); rows in date order, not to scale. Markers give the lane: legacy-split SSL and recipe work,
registered E1 openings, benchmark protocol. Legacy-split rows are not comparable with Chapters~\ref{ch:methods} and~\ref{ch:e1_result};
test50 values are in-distribution, known leakage, context only; opening \#8: C2 recipe averaged over seeds 5 to 9; OP-A; 101 of 106
cells CSIR.}

\section{Conditions of every number in this chapter}\label{sec:legacy-conditions}

All work of the period used the legacy split (478 train, 54 validation, 50 test cells), trained on label set v1 unless stated
(Section~\ref{sec:data-labels}), and took its primary contrasts on the 50 test50 cells. E1 was read without registration during the P-series (Table~\ref{tab:T21}),
and from 1 Nov some protocols registered E1 co-primaries. Arms were compared under a locked rule fixed before each run: a difference passes if it exceeds
$\max(0.01,\, 2\sqrt{(s_A^2 + s_B^2)/3})$ on any primary with no primary worse, where $s_A$ and $s_B$ are the standard deviations over
seeds; a difference is called \emph{robust} if its 95\% image-bootstrap interval excludes 0 and its sign survives removing the top 3
cells. Outcome words in Table~\ref{tab:T21} (supported, n.d., WORSE) are the words of that rule.

These numbers can show the direction of a paired contrast between two arms trained on the same split and scored on the same cells
under the rule in force. They cannot show four things, and we do not use them for any of them.
\begin{enumerate}
  \item \textbf{Comparable levels.} Split, label set, recipe and rule differ from Chapters~\ref{ch:methods} and~\ref{ch:e1_result}; no
    level here can be set beside a level there.
  \item \textbf{Generalisation beyond the scored cells.} test50 overlaps the training pool: 8 test cells are the same physical cell as
    a training cell and 7 more are neighbouring crops (Section~\ref{sec:data-split}). Every test50 number is in-distribution with
    known leakage, context only.
  \item \textbf{A blind result.} test50 was scored in every run of the period, and E1 was read many times (Section~\ref{sec:legacy-taught}).
    No P-series model was selected on test50 (final weights, validation-only choices), but both sets were seen so often that later
    choices cannot be called blind to them.
  \item \textbf{Small effects.} With 50 test cells, the smallest detectable crack IoU change was about 0.012 at five against five
    seeds, and inactive results were decided by 3 cells. The rule itself was lenient: about 2 standard errors, roughly a one-sided
    $p \approx 0.06$ at 3 seeds, where a two-sided 95\% test needs 2.8. ``Supported'' under this rule is not ``confirmed'' under the
    protocol.
\end{enumerate}

\section{The semi-supervised period (10 to 13 Oct)}\label{sec:legacy-ssl}

The project began as a study of semi-supervised learning (SSL): pseudo-label methods in the Mean Teacher family
\citep{tarvainen2017meanteacher} that use unlabelled EL cells beside a small labelled fraction. At first the supervised
baseline beat every SSL arm. An audit of the design, not of the decision rule, found five causes:
\begin{enumerate}
  \item \textbf{Regime.} The runs used 5 to 10\% labels, where inactive occurs in only 2 or 4 labelled cells and the teacher is weakest.
  \item \textbf{A built-in handicap.} Every SSL arm until P13 carried one: first a stronger labelled view, then batch-norm statistics
    shared between labelled and unlabelled images. With shared statistics the loss-off control already lost crack IoU against plain
    supervision ($-0.017$ $[-0.024, -0.010]$ at 5\% labels) before the
    unlabelled loss did anything.
  \item \textbf{Seeds.} Three seeds were used for effects of 0.01 to 0.02.
  \item \textbf{Tuning asymmetry.} The supervised recipe was chosen by a four-step ladder on validation; the SSL hyperparameters were
    fixed in advance and never tuned.
  \item \textbf{Data.} The in-domain unlabelled pool was 239 to 455 cells; the external pool came from two sources, half of it low in contrast.
\end{enumerate}
The audit also found the perturbation weak: the unlabelled loss was about 10 $\times$ smaller than the supervised loss, and CutMix
\citep{yun2019cutmix} raised the ratio to 0.77.

With separate batch-norm statistics and CutMix, the P13 primary at 5\% labels (five against five seeds) passed: inactive IoU
$+0.117$ $[+0.042, +0.189]$ and crack relF1@2 $+0.033$ $[+0.006, +0.055]$, with crack IoU $+0.014$ n.d.; a secondary arm with hard
pseudo-labels gained $+0.029$ $[+0.014, +0.041]$ crack IoU (two of its five seeds added after a three-seed result). The
separate-statistics recipe (with CutMix dropped by a validation pick) was n.d. at 20\%, 50\% and 100\% labels (P16). With the
external pool of 11,938 unlabelled CSIR and ARTS cells, the pseudo-label loss did not help against its loss-off control at 5\%
labels: crack IoU $-0.012$ $[-0.026, -0.001]$ (robust over cells, n.d. by seeds) and relF1@2 $-0.027$, worse by the rule (P10, still
with shared batch-norm statistics). Encoder pretraining on the related UCF-EL dataset was worse on all three primaries at 10\% labels (P15, 3 seeds, under the
house fine-tune that Section~\ref{sec:legacy-pf} shows to erase the deep encoder), and an
adapted region-routing SSL method \citep{lei2025alhvr} was WORSE on all three primaries in all four cell sets (AL1).

Our reading: SSL helped this model at 5\% labels and no gain was detectable from 20\% labels up, for the recipe that was run. This
closes only those conditions (shared or separate batch-norm statistics, a ResNet-34 U-Net++, an unlabelled pool of these two
sources, three to five seeds). It does not show that SSL fails on EL images. The work then turned to models trained on all labels (from 24 Oct),
which is what the goal of this thesis concerns (Chapter~\ref{ch:intro}), and the SSL framing was later dropped.

\section{Labels moved the numbers more than the SSL losses}\label{sec:legacy-labels}

The clearest pattern of the period was that the label version moved crack scores more than any SSL effect did. Re-scoring old models on
the revised labels raised crack IoU by $+0.05$ to $+0.07$ (P3, exploratory). Training on revised rather than old labels at 10\%
labels raised it by $+0.082$ $[+0.046, +0.129]$ on test50 (in-distribution, known leakage, context only), and by $+0.030$ $[+0.012, +0.048]$ on the E1 cells (P14). A loss that tolerates
offsets of the human labels recovered none of the label-noise gap ($-0.004$). On the supervised side, label-safe augmentation (flips,
quarter turns, brightness and contrast only) helped (crack AP $+0.020$, inactive AP $+0.028$; P6), and going from 10\% to 100\% of
the labels raised crack IoU by $+0.075$ and inactive IoU by $+0.135$. Choosing which cells to label also mattered: 12 crack cells
added to a 5\% base gained $+0.030$ $[+0.020, +0.042]$ relF1@2 over 12 random cells, but the crack cells were chosen with the labels
known (oracle-selected). The audit of Chapter~\ref{ch:data} later measured the same effect on the pool: the label release moves PF's
relF1@2 by more than the final recipe gains over PF on CV5 (selection-biased; Section~\ref{sec:data-labels}).

\section{The encoder-erasing fine-tune and PF (30 Oct to 1 Nov)}\label{sec:legacy-pf}

The house recipe trains U-Net++ \citep{zhou2018unetpp} with an ImageNet ResNet-34 encoder by Adam at learning rate $10^{-3}$ with
coupled weight decay $10^{-4}$ for 20,000 steps. A direct measurement of the weights showed that this fine-tune erases the deep
encoder: the per-tensor cosine between the final and the initial convolution weights is 0.005 to 0.018 in layer3 and $-0.001$ to
0.001 in layer4, for ImageNet and for masked-image-modelling initialisations, at 5\% and 100\% labels. The direction is lost within
the first 2,000 steps (layer3 and layer4 at 0.070 and 0.003). This finding changed the reading of earlier nulls. Pretraining on the
unlabelled pool by masked image modelling or decoder denoising \citep{brempong2022ddep} had been n.d. at 5\% labels (MIM5, DDeP5),
but under this fine-tune the masked-image run could only test the transfer of the first encoder layers.

The encoder-preserving fine-tune (PF) freezes the encoder for 1,000 steps, then ramps it to 0.1 times the learning
rate over the next 1,000 steps, with no encoder weight decay; layer3 and layer4 stay at cosine 0.996 and 0.997 to their initial weights after 2,000 steps. Under PF,
on new seeds 5 to 9 (E1 opening \#1), E1 inactive IoU rose from 0.0481 to 0.2371 ($+0.1890$ against the rule's threshold 0.0433),
replicating an earlier result on the same 106 E1 cells. On test50 (in-distribution, known leakage, context only) the crack relF1@2
gain was $+0.0117$ on all 50 cells and $+0.0120$ on the clean 49, a pass by 0.0017 and 0.0020 that rests on a few cells; E1 crack
recall was n.d. As five-member ensembles (E1 opening \#4, not independent of \#1), PF raised E1 inactive IoU by $+0.2141$
$[+0.0303, +0.3820]$, with test50 n.d. (context only). A module-grouped five-fold read on the pool (B6, 2 seeds) found no detectable in-distribution inactive difference
($-0.0144$ $[-0.0977, +0.0570]$) and excluded an in-distribution gain as large as the E1 gain. ResNet-50 under PF was worse than
ResNet-34 on test50 inactive IoU ($-0.0289$, threshold 0.0251; context only, as above).

PF became the baseline of the protocol and of the final comparison (C1 in Chapter~\ref{ch:e1_result}) because of these E1 reads.
The choice of the baseline was therefore informed by E1, which is one of the prior looks that Chapter~\ref{ch:e1_result} lists.

\section{Other arms of the period}\label{sec:legacy-arms}

\begin{itemize}
  \item \textbf{Skeleton Recall} \citep{kirchhoff2024skeletonrecall} at weight 1 (E1 opening \#2) raised E1 crack relF1@2 by $+0.0590$
    (threshold 0.0392) but lowered test50 crack IoU by $-0.0136$ (threshold 0.0100; in-distribution, known leakage, context
    only), and failed its pre-registered rule. At weight
    0.1 (opening \#5) it passed one E1 co-primary narrowly ($+0.0364$ against 0.0351, about 10\% family-wise chance of a false pass
    over 7 co-primaries); as an ensemble it failed two guards (opening \#6, not independent of \#5); on fresh seeds 10 to 14 (opening
    \#7, 14 Nov) the E1 difference appeared again, but replication at recipe level was not shown (seed-aware interval
    $[-0.0015, +0.0590]$). Soft clDice failed; scuff as a class of its own was not learned (train scuff IoU 0.000 in all 5 seeds).
  \item \textbf{Synthetic half-cut cells} (opening \#3) were n.d. on test50 (context only) and E1 (E1 relF1@2 $+0.0135$ against 0.0429).
  \item \textbf{A foundation encoder.} DINOv3 ConvNeXt-T \citep{simeoni2025dinov3} against ImageNet ConvNeXt-T under PF (2 seeds,
    no test50, no E1; crack AP on 51 E-SCDD full cells) gave $-0.0243$ with a lower bound of $-0.0463$, where continuing needed a
    positive bound: stopped, with ``no sign of a gain'', not ``worse''.
  \item \textbf{Single-model changes at 100\% labels} (Track M; test50, in-distribution, known leakage, context only): Dice on 12 classes was worse on crack ($-0.0106$), ResNet-50 worse on
    inactive ($-0.0251$), 107 extra E-SCDD cells and twice the input resolution n.d.; under PF, masked-image pretraining and crack
    copy-paste were n.d. at 100\% labels.
  \item \textbf{The five-seed ensemble} with a crack bias beat its single members by $+0.018$ $[+0.013, +0.025]$ crack IoU on the clean
    49 test50 cells, under a locked protocol with validation-only choices (in-distribution, known leakage, context only). This was the
    one clear win of the period.
\end{itemize}

\section{What the period taught}\label{sec:legacy-taught}

Four lessons of this period became rules of the protocol (Chapter~\ref{ch:protocol}).
\begin{enumerate}
  \item \textbf{The test set was used up.} test50 was scored in every experiment before 24 Oct without being counted, then opened 18
    counted times up to 14 Nov. It overlaps the training pool (Section~\ref{sec:data-split}), 3 of its cells hold 59\% of its inactive
    pixels, and one very noisy cell wrecked about one run in six unless excluded. The protocol therefore develops on module-grouped
    cross-validation of the pool and keeps test50 for continuity only.
  \item \textbf{E1 was read too often.} Seven registered openings on the same 106 cells (1 to 14 Nov; Table~\ref{tab:T17}), three of
    them looks at one contrast and two not independent of an earlier one, came on top of reads before the 31 Oct policy that were not
    counted. The protocol allows one registered opening with candidates fixed before any number.
  \item \textbf{The instrument was too weak.} With 50 cells and a lenient rule, a crack IoU change below about 0.012 could not be seen,
    inactive was decided by 3 cells, and seven co-primaries gave a false pass about one chance in ten. The protocol uses 532 pool cells out of fold,
    five seeds, a margin from an A/A split and one primary per arm.
  \item \textbf{Strict crack IoU and the label version dominated.} Label changes moved crack IoU by more than any SSL effect did, and
    70\% of the best model's crack IoU loss lay within 2 px of the label. The protocol lets relF1@2 decide and reports strict IoU with its
    convention band (Section~\ref{sec:data-f1}).
\end{enumerate}
What survives from the period are mechanisms and directions, each measured under its own conditions: label quality matters, the
house fine-tune erases the deep encoder, PF raises inactive IoU on E1, and an ensemble helps on test50 (context only). Chapters~\ref{ch:methods}
and~\ref{ch:e1_result} re-test what mattered under the protocol.
